\bookchapter{Kapitel 26: Aktuelle Fronten der Smartphone-Technologie}

\chapterintro{Nachdem wir das „ digitale Erbe“ alter Handys besprochen haben, wenden wir uns der Gegenwart zu. 2024-2026 ist eine Phase geballter Smartphone-Innovationen: KI ist von der Cloud in die Hosentasche gewandert, die Satellitenkommunikation vom Flaggschiff ins Mittelsegment gewandert, faltbare Displays vom „Experimentieren” zum „taeglichen Gebrauch“ geworden, und auch die Batterietechnologie erlebt einen seit Jahren nicht mehr dagewesenen Materialdurchbruch. Dieses Kapitel beleuchtet diese aktuellen Technikfronten und hilft dir einzuschaetzen, ob „es Zeit fuer ein neues Handy ist” und „welches“.}

\booksection{I. KI-Smartphone: Marketing-Schlagwort oder echte Revolution?}

Seit 2024 bewerben fast alle Android-Hersteller ihre Flaggschiffe als „KI-Handy“; auch Apple hat mit iOS 18 Apple Intelligence eingefuehrt. Aber was genau macht ein Handy zum „KI-Handy”? In der Branche haben sich drei Beurteilungskriterien herausgebildet:

\begin{itemize}
  \item \textbf{Warum ist die Speichergeschwindigkeit wichtig?}: Das Handy kann lokal ein Sprachmodell mit ca. 7B Parametern ausfuehren (wie Qwen, Llama, DeepSeek in leichter Version), ohne Cloud-Abhaengigkeit Aufgaben wie Zusammenfassen, Uebersetzen und Textgenerierung erledigen. Testmethode: Wenn die KI-Funktion im Flugmodus noch funktioniert, ist es ein On-Device-Modell.
  \item \textbf{NPU-Rechenleistung (TOPS)}: Die Rechenleistung der NPU (Neural Processing Unit) wird in TOPS (Billionen Operationen pro Sekunde) angegeben. 2024 erreichten Flaggschiff-Chips (wie Snapdragon 8 Gen3, Dimensity 9300) 40-70 TOPS; dieNeue Generation Chips 2025-2026 uebertreffen allgemein 100 TOPS – das ist die Hardware-Huerde fuer das lokale Ausfuehren grosser Modelle.
  \item \textbf{KI-Agent-Szenarien}: Nicht nur „Fragen und Antworten", sondern auch die faehige Ausfuehrung von Aufgabenketten ueber mehrere Apps hinweg – z. B. „Sortiere die Fotos von gestern in ein Album und teile sie mit Zhang San"; das System ruft automatisch Album, Datei-Manager, WeChat und andere Apps auf. Das ist das Kernmerkmal, das ein KI-Handy von einem „Handy mit KI-Assistent" unterscheidet.
\end{itemize}

\booksection{II. Vergleich der KI-Strategien der Hersteller: Wer macht es ernst, wer klebt nur ein Etikett auf?}

Die KI-Wege der Hersteller unterscheiden sich deutlich; beim Kauf lohnt sich ein genauerer Blick:

\begin{itemize}
  \item \textbf{Apple Apple Intelligence (iOS 18+)}: Schwerpunkt auf Integration in System-Apps – E-Mail-Zusammenfassungen, Benachrichtigungsprioritaet, Schreiboptimierung, KI-Bildbearbeitung in der Galerie. Primaer On-Device-Modell; komplexe Aufgaben werden an Private Cloud Compute ausgelagert. Vorteile: Datenschutz und Systemfluessigkeit. Nachteil: Chinesischsprachige Funktionen kamen spaeter; einige Features erst ab 2025.
  \item \textbf{Google Gemini Nano}: In Pixel 8/9-Serie integriert; unterstuetzt „Circle to Search", Magic Compose, Audio-Zusammenfassungen usw. Gemini Nano ist ein echtes On-Device-Kleinmodell, quelloffen und freundlich zur Open-Source-Community – die Referenzimplementierung fuer Android-KI-Handys.
  \item \textbf{: Am Sehenswürdigkeit ein fantastisches Selfie geschossen, aber der Hintergrund ist voller anderer Touristen? Mit „Magie-Entfernung" die Passanten umranden – das Smartphone erkennt sie automatisch, „malt" sie weg und füllt den Hintergrund intelligent auf. Das Ergebnis ist oft verblüffend gut.}: Anbindung an HarmonyOS XiaoYi; unterstuetzt On-Device-Textgenerierung, KI-Bildbearbeitung, KI-Entfernung (Passanten wegnehmen) und Anrufzusammenfassungen. Mit HarmonyOS NEXT 5.0 weitere systemweite Integration; KI-Faehigkeiten verschmelzen tief mit dem System.
  \item \textbf{Xiaomi MiLM (Mi Large Model)}: In HyperOS integriert; unterstuetzt KI-Portraets, KI-Zusammenfassungen, KI-Malwerkzeug und Super-Aufloesung. Xiaomi verfolgt eine „On-Device + Cloud"-Strategie: leichtes Modell lokal, schwere Aufgaben in der Cloud.
  \item \textbf{OPPO (Andes-Modell) und vivo (BlueLM-Modell)}: In ColorOS 15 und OriginOS 5 tief integriert; Schwerpunkte: KI-Entfernung, KI-Dokumentzusammenfassung, KI-Bilderkennungs-Q\& A und natuerlichsprachige Fotosuche („Finde die Fotos vom letzten Jahr am Meer").
\end{itemize}

\textbf{Kaufempfehlung}: Wenn du regelmaessig KI fuer Arbeitsdokumente und Informationsorganisation brauchst, waehle ein Geraet mit starker On-Device-Leistung und systemintegrierten KI-Funktionen. Fuer gelegentliche KI-Bildbearbeitung reicht ein Mittelklasse-Handy.

\booksection{III. Satellitenkommunikation: vom Flaggschiff-„Gag“ zum Mittelklasse-Standard}

Die Satellitenkommunikation ist eine der substanzellsten Technologien, die in den letzten zwei Jahren von der Oberklasse ins Mittelsegment gewandert sind. Als Huawei Mate 60 und iPhone 14 sie 2022 ein fuehrten, war sie noch Flaggschiff-exklusiv; 2025 ist sie bereits im 2000-3000-RMB-Segment angekommen.

\begin{itemize}
  \item \textbf{Beidou-Kurzmitteilung im Mittelsegment}: Die Kurzmitteilungsfunktion des Beidou-Satelliten (SMS mit bis zu 1000 chinesischen Zeichen + Position) ist in Mittelklasse-Modellen von Redmi, realme, iQOO usw. verfuegbar – ohne Grundgebuehr, kostenlos. Geeignet fuer Wanderungen und Roadtrips.
  \item \textbf{Tiantong-Satellitentelefonie}: Huawei Mate 60/70-Serie, Honor Magic6-Serie usw. unterstuetzen die Tiantong-1-Satellitentelefonie; Satellitengespraeche anrufen/empfangen und Satelliten-Nachrichten senden. Ein Satellitentarif des Anbieters ist erforderlich (ca. 10-200 RMB/Monat); fuer Nutzer, die oft in unbewohnte Gebiete gehen.
  \item \textbf{iPhone-Satellitennachricht (SOS)}: Ab iPhone 14 wird Satelliten-SOS unterstuetzt; Textnachrichten koennen an Rettungsdienste gesendet und die Position geteilt werden. 2 Jahre kostenlos; in einigen Regionen verlaengerbar.
  \item \textbf{: Ihre Stärke sind Playlists, Community und Kommentare. Sie möchten „Musik für Regentage" oder „entspannende Instrumentalmusik zum Lernen" hören? Eine Suche genügt, und Sie finden von anderen kuratierte Playlists. Manchmal sind die Kommentare besser als der Song selbst. Nachteil: Bei einigen Songs fehlen die Rechte – sie werden während des Hörens ausgegraut.}: Stadtbewohner brauchen es kaum; fuer Outdoor-Enthusiasten und Geschaeftsreisende in abgelegenen Gebieten wird es dringend empfohlen – in Notfaellen kann es Leben retten.
\end{itemize}

\booksection{IV. Reife des faltbaren Displays: Bereit als Hauptgeraet?}

Faltbare Displays sind vom „Tech-Experiment“ zum „tauglichen Hauptgeraet” geworden. Drei Kriterien fuer die Reife: Hardware-Haltbarkeit, App-Anpassung und Interaktionsoekosystem.

\begin{itemize}
  \item \textbf{Hardware-Haltbarkeit}: Scharnierlebensdauer ueber 500.000 Oeffnungen (bei 100-mal taeglich 13 Jahre); Tropfen-Scharnier reduziert den Knick; IPX8-Wasserschutz ist Standard. UTG-Glas + Schutzfolie machen den Innenbildschirm weniger empfindlich.
  \item \textbf{App-Anpassung}: Die wichtigsten Apps (WeChat, Taobao, Douyin, Bilibili, Office, WPS) sind auf den grossen Bildschirm optimiert; aber ca. 30 \% der Nischen-Apps sind noch nicht angepasst und erscheinen gestreckt oder mit schwarzen Raendern. Vor dem Kauf die haeufigsten Apps ausprobieren.
  \item \textbf{Halb-geklappte Interaktion}: Bei 90°-120°-Klappung unterstuetzen einige Apps einen Modus „oben Inhalt, unten Steuerung" – geeignet fuer Videoanrufe, Selfies und Video-Wiedergabe. Die Selfie-Halbklapphaltung ist einzigartig fuer Falthandys.
  \item \textbf{Split-Screen und Multitasking}: Auf dem grossen Bildschirm koennen zwei Apps im Split-Screen oder drei Apps in kleinen Fenstern gleichzeitig laufen. Die Split-Screen-Logik von HyperOS, ColorOS und OriginOS ist ausgereift; iOS (im iPadOS-Stil) bietet auf dem iPhone noch keinen Split-Screen.
  \item \textbf{Preis}: Querfalz (Buch-Stil) von ueber 10.000 RMB auf 5.000-7.000 RMB gesunken; Laengsfalz (Klapp-Handy) bereits bei 4.000-5.000 RMB. 2025 bringen einige Hersteller Laengsfalz-Modelle um 3.000 RMB – faltbare Displays erreichen den Massenmarkt.
\end{itemize}

\textbf{Kaufempfehlung}: Fuer Vielleser, Dokumentbearbeiter und Video-Nutzer lohnt sich ein Querfalz-Display; wer Wert auf Tragbarkeit und soziale Attraktivitaet legt, ist mit einem Laengsfalz besser beraten. Vor der Entscheidung mindestens 30 Minuten im Laden ausprobieren.

\booksection{V. 6G und Wi-Fi 7: Tempo und Grenzen der naechsten Verbindungsgeneration}

6G ist noch nicht wirklich angekommen, aber die Technologie ist bereits auf dem Weg:

\begin{itemize}
  \item \textbf{6G-Fortschritt}: 3GPP plant, den 6G-Standard bis 2027 fertigzustellen; kommerzielle Nutzung ca. 2030. Die theoretische Spitzenrate von 6G koennte bei 1 Tbps liegen (50x so schnell wie 5G), die Latenz bei 0,1 ms. „Integrated Sensing and Communication" (Kommunikation + Wahrnehmung) wird unterstuetzt; das Handy koennte sogar als grob aufloesendes Radar dienen. Die Hersteller befinden sich noch im Prototyp-Stadium; Consumer-Klasse Produkte sind noch fern.
  \item \textbf{Wi-Fi 7(802.11be)}: Ab 2024 Standard bei Flaggschiffen; theoretische Spitzenrate bis 46 Gbps (4x Wi-Fi 6); unterstuetzt 320 MHz Bandbreite, 4K-QAM und MLO (Multi-Link-Uebertragung). Gemessen an einem Wi-Fi-7-Router erreichen Flaggschiff-Handys 2-5 Gbps – 2-3x schneller als Wi-Fi 6E. Wer ein Wi-Fi-7-Handy kauft, sollte auch einen Wi-Fi-7-Router haben, um die volle Leistung zu erhalten.
  \item \textbf{Bluetooth 5.4 und LE Audio}: 2024-2026 unterstuetzen neue Handys allgemein Bluetooth 5.4 und LE Audio (Low Energy Audio); zwei Kopfhoerer-Paare koennen gleichzeitig dasselbe Lied hoeren. Auracast-Broadcast-Audio ermoeglicht es vielen Personen an oeffentlichen Orten, dieselbe Tonquelle synchron zu hoeren.
\end{itemize}

\booksection{VI. Silizium-Kohlenstoff-Anode: Die Akkulaufzeit erfaehrt endlich einen Materialdurchbruch}

Die Akkulaufzeit von Handys wurde jahrelang als „taegliches Laden“ kritisiert; seit 2024 hat die Verbreitung der Silizium-Kohlenstoff-Anoden-Technologie eine spuerbare Verbesserung gebracht.

\begin{itemize}
  \item \textbf{Technisches Prinzip}: Die theoretische Kapazitaet herkoemmlicher Graphitanoden liegt bei ca. 372 mAh/g; die von Silizium bei 4.200 mAh/g (ca. 10x). Beim Laden/Entladen dehnt sich Silizium jedoch um 300 \% aus und neigt zur Pulverisierung. Die Silizium-Kohlenstoff-Anode kontrolliert diese Ausdehnung durch eine Nano-Silizium-Kohlenstoff-Verbundstruktur in akzeptablem Rahmen und erhoeht die Energiedichte um 15-30 \%.
  \item \textbf{Tatsächliche Leistung}: Flaggschiffe 2024-2025 mit Silizium-Kohlenstoff-Anode (wie Xiaomi 14 Ultra, Honor Magic 6 Ultimate, vivo X100 Pro) bieten in 7-8 mm duennem Gehaeuse 5.000-6.000 mAh – ca. 800-1.000 mAh mehr als die Vorgaenger. Bei intensiver Nutzung reicht der Akku ca. 1,5 Tage.
  \item \textbf{100-Watt-Schnellladen als Ergaenzung}: Die Silizium-Kohlenstoff-Anode ist gut mit Schnellladen kompatibel; 100 W/120 W in 17 Minuten voll ist keine Neuigkeit mehr. 2025 unterstuetzen einige Modelle 150-240 W; 80 \% in 10 Minuten.
  \item \textbf{Kaufempfehlung}: Als intensiver Nutzer waehle vorzugsweise ein Modell ab 2024 mit „Silizium-Kohlenstoff-Anode + 100-W-Schnellladen"; die Reichweitenangst wird deutlich geringer.
\end{itemize}

\booksection{VII. Kapitelueberblick: Diese Fronten in einem Satz merken}

\begin{itemize}
  \item Drei KI-Kriterien: On-Device-Modell, NPU-TOPS und KI-Agent mit App-uebergreifenden Szenarien.
  \item Satellitenkommunikation: Beidou-Kurzmitteilung im Mittelsegment; Tiantong-Telefonie noch Flaggschiff-exklusiv.
  \item Faltbare Displays: Hardware ausgereift, App-Anpassung ca. 70 \%, Halbklapp- und Split-Screen sind einzigartig.
  \item 6G ist unterwegs, Wi-Fi 7 ist angekommen, Bluetooth LE Audio beginnt sich durchzusetzen.
  \item Silizium-Kohlenstoff-Anode + 100-W-Schnellladen: Doppeltdurchbruch bei Laufzeit und Laden.
\end{itemize}

\begin{tcbpractice}{Ausprobieren: Erstelle eine „Technik-Anforderungsliste" fuer dein naechstes Handy}
Nimm ein Blatt Papier und ordne deine Wuensche fuer KI, Satellit, faltbares Display, Wi-Fi 7, Schnellladen und Akkukapazitaet in drei Stufen: „unbedingt / Bonus / nicht noetig“. So kannst du beim naechsten Besuch im Laden schnell beurteilen, ob ein vom Verkaeufer angepriesenes „Verkaufsargument” wirklich dein Beduerfnis trifft. Beispiel:

\begin{itemize}
  \item Unbedingt: Akku ueber 5.000 mAh, Schnellladen ueber 80 W, NFC, IP68-Wasserschutz.
  \item Bonus: On-Device-KI-Zusammenfassung, Wi-Fi 7, Beidou-Kurzmitteilung.
  \item Nicht noetig: faltbares Display (ich mag Bar-Handys), Satellitentelefonie (ich reise nicht in die Wildnis), 240-W-Extrem-Schnellladen.
\end{itemize}

\end{tcbpractice}

\begin{tcbsummary}{Kapitelzusammenfassung}
\begin{itemize}
  \item Drei Kriterien fuer KI-Handys: On-Device-Modell, NPU-TOPS, KI-Agent-Szenario
  \item KI-Strategien: Apple – Datenschutz und Integration; Google – On-Device und Open Source; Huawei/Xiaomi/OPPO/vivo – je eigene Staerken
  \item Satellitenkommunikation: Beidou-Kurzmitteilung im Mittelsegment; Tiantong-Telefonie noch Flaggschiff-exklusiv
  \item Faltbare Displays: Hardware ausgereift, ca. 70 \% App-Anpassung; Halbklapp- und Split-Screen einzigartig
  \item 6G-Standard 2027, Wi-Fi 7 verbreitet, Bluetooth LE Audio im Kommen
  \item Silizium-Kohlenstoff-Anode macht 5.000-6.000 mAh zum Flaggschiff-Standard; Reichweitenangst deutlich verringert
\end{itemize}

\end{tcbsummary}
